\section{Social Engineering}\label{sec:SocialEngineering}
In questa sezione analizzeremo la \emph{Social Engineering} dietro al malware trattato.

\subsection{Curiosità ed emozioni}
Ricevere un email da una persona con oggetto "ILOVEYOU" e corpo "Kindly check the attached LOVELETTER coming from me." fa crescere la curiosità nell'animo della vittima. 
Il messaggio particolarmente breve e generico contribuiva a far credere alla vittima che l'email fosse legittima qualunque fosse il mittente.

Si consideri anche che l'esca emotiva utilizzata è proprio l'amore; 
infatti tra tutte le emozioni venne scelta una tra le più forti, a parimerito pochissime altre emozioni, 
come ad esempio la paura che tratteremo nel seguito del capitolo~\ref{subsec:fear}.

Si provi ad immaginare ora la sopresa di ricevere attenzioni amorose inaspettate da un/una collega, un amico/un'amica oppure da un parente
\begin{joke} 
  in Alabama.
\end{joke}
Questo sarebbe un fortissimo incentivo a controllare il contenuto, a causa della natura curiosa dell'essere umano. 

\subsection{Falsa sicurezza}
Di fatti le email venivano ricevute da persone, precedentemente infettate, che avevano il nostro indirizzo email sul proprio "Outlook address book"; 
si suppone dunque che i mittenti avessero un certo grado di confidenza con la vittima.
Più il mittente è fidato più aumentavano le possibilità del file di essere aperto.

\subsection{Errore umano nella lettura del file}
Chiunque conosce l'estensione ".txt" e praticamente tutti la considerano non malevola.
Il malware utilizzava una feature di Windows attiva di default che nascondeva l'estensione del file se riconosciuta dal sistema.

Accadeva dunque che un file nominato "LOVE-LETTER-FOR-YOU.TXT.vbs" venisse visualizzato come "LOVE-LETTER-FOR-YOU.TXT", rafforzando dunque l'idea che si trattasse di una innoqua lettara d'amore.

\subsection{Paura}\label{subsec:fear} %#TODO: se annullata cancellare citazione in Curiosità ed emozioni
